\section{Experimental Plan}
\label{sec:experimental-plan}

% =========================================================================
% Pre-registered evaluation plan. The function of this section is to make
% the eventual results section harder to game and easier to evaluate: the
% baseline grid, the compute axis, the headline figure, and the
% falsification criterion are all fixed before any training begins.
% =========================================================================

This section pre-registers the empirical study that the framework of
Sections~\ref{sec:method-state}--\ref{sec:method-inference} is intended
to support. The plan is fixed before training; subsequent empirical work
either runs against this plan or explicitly documents departures from
it. The intent is to draw a clean line between the framework
contribution (this paper) and the empirical contribution (a follow-up
study), and to make the eventual headline number harder to confuse with
post-hoc tuning. The structure of the plan mirrors the methodological
commitments of Section~\ref{sec:method}: the agent is an approximate
$Q$-learning system over a verifier-anchored context MDP, MCTS is a
policy improvement operator used during training and a test-time-compute
knob at inference, and the falsifiable claim concerns the
compute-versus-accuracy frontier rather than any single accuracy number.

\subsection{Benchmarks}
\label{sec:exp-benchmarks}

We evaluate on a progression of verifiable-math benchmarks ordered from
easiest to hardest: GSM8K \citep{cobbe2021gsm8k}, the MATH dataset
\citep{hendrycks2021math} split into levels 1 through 3 and then levels
4 through 5, AIME 2024 and 2025, and Lean Mathlib as a deferred option
for formal-proof evaluation. The progression matters operationally.
Approximate $Q$-learning over the context MDP of
Section~\ref{sec:method-state} needs at least some nonzero terminal
reward to bootstrap: if the base model scores near zero on a benchmark,
the verifier returns a constant signal across visited trajectories, the
Bellman target collapses to a point estimate, and there is no gradient
on which the heads or the MCTS-distilled policy can sharpen. Easier
benchmarks supply the early signal that drives initial policy
improvement; the harder benchmarks then test whether the trained
$\Qpsi$ generalizes to problems the base model could not solve on its
own. Starting at the hard end is operationally tempting (it is where
the headline gap to frontier models is widest) but methodologically
risky: a flat-zero learning signal does not falsify the framework, it
simply provides no evidence either way.

Each benchmark is paired with an explicit, non-LLM verifier whose
implementation is part of the pre-registered plan. For GSM8K we regex
the boxed numeric answer from the committed output, canonicalize it
with SymPy, and compare to the ground truth via SymPy's equality. For
MATH we use the same regex and SymPy stack, extended to handle
symbolic expressions (fractions, radicals, algebraic forms); the
verifier accepts any expression that SymPy reports as equal to the
ground-truth answer under canonicalization. For AIME the answer is a
nonnegative integer in $[0, 999]$, so the verifier is integer parsing
plus equality. For Lean Mathlib the verifier is the kernel: a tactic
proof is accepted iff Lean's type checker reports a closed proof of the
stated goal. The Lean target is deferred because it requires a
substantially larger engineering investment around tactic
serialization, error-message ingestion, and dependency management; the
framework is constructed to admit it without redesign, but the phase-1
empirical study targets the SymPy-checkable benchmarks first. The
verifier interface is the same callable referenced by the
\textsc{Verify} action of Section~\ref{sec:method-actions}: the
agent can spend intermediate verifier calls during deliberation and
also be scored by the same checker at terminal commit.

\subsection{Base Models}
\label{sec:exp-models}

The base language model is the frozen substrate of the architecture in
Section~\ref{sec:method-architecture}, and its identity has to support
the central comparison framing: any measured improvement during
training is attributable to the heads, the adapter, and the search
procedure, not to capability gains from prior reasoning-specific
post-training. We therefore restrict the substrate to small base models
that have not been instruction-tuned, RLHF-tuned, or specialized for
reasoning. The primary candidate is Qwen2.5-1.5B-Base
\citep{yang2024qwen25}; secondary candidates are Llama-3.2-1B-Base
\citep{dubey2024llama3} and SmolLM2-1.7B \citep{allal2024smollm2}. The
heads and adapter are trained from scratch; the LLM stays frozen
throughout phase 1, in line with the architectural choice of
Section~\ref{sec:method-architecture}.

We deliberately exclude Instruct, Chat, RLHF, o1-style, and R1-style
variants \citep{openai2024o1, guo2025deepseekr1} from the substrate
pool. Such variants already encode reasoning behavior obtained from
methods we want to compare against, not build on. Using them would
conflate two sources of improvement: gains from approximate
$Q$-learning with MCTS-supplied targets (the framework's contribution)
and gains from the prior post-training pipeline (someone else's
contribution). A base-only substrate makes the attribution clean. The
larger reasoning model that appears in the baseline grid as B5 (below)
is not the substrate; it serves only as a capability-ceiling reference
point for the compute-versus-accuracy frontier, and its identity is
allowed to differ across benchmarks if a stronger published reasoning
model exists for one but not another.

\subsection{Pre-Registered Baseline Grid}
\label{sec:exp-baselines}

The baseline grid in Table~\ref{tab:baselines} is fixed before training
begins and is not extended post-hoc. Pre-registration matters because
the natural failure mode of test-time-compute experiments is exactly
the kind of selective baseline reporting that the grid is designed to
prevent: a method whose primary comparison is to a weak control wins
trivially, while a method whose primary comparison is to a strong
control loses trivially. Fixing the grid in advance, with explicit
roles for each baseline, removes that degree of freedom.

\begin{table}[t]
\centering
\caption{Pre-registered baseline grid.}
\label{tab:baselines}
\small
\begin{tabular}{@{}llp{0.5\linewidth}@{}}
\toprule
ID & Method & Description \\
\midrule
B0 & Base greedy        & Base LLM with greedy decoding; no deliberation. \\
B1 & Best-of-$N$        & Base LLM with $N$-sample best-of-$N$, reported as a
                          swept curve over $N$ and FLOP-matched per problem to
                          B4 at the headline comparison point. \\
B2 & Policy $B{=}0$     & Frozen LLM + learned policy/$Q$ heads; fast mode,
                          no search. \\
B3 & MCTS $B{=}N_{\mathrm{small}}$ & Same as B2 plus light MCTS at
                          inference. \\
B4 & MCTS $B{=}N_{\mathrm{large}}$ & Same as B2 plus heavy MCTS at
                          inference. \\
B5 & Larger model       & Larger reasoning model, greedy decode, as a
                          capability ceiling. \\
\bottomrule
\end{tabular}
\end{table}

The roles assigned to the six baselines are deliberately
non-overlapping. B0 is the raw capability floor: greedy decoding from
the base model with no deliberation scaffold of any kind. The B0 score
is the number that the rest of the grid is trying to beat. B1 is the
cheap-deliberation control: best-of-$N$ sampling from the base model,
the canonical repeated-sampling baseline whose inference-compute scaling
\citet{brown2024monkeys} characterize, reported here as a swept curve over
$N$. The single FLOP-matched comparison
point used in the falsification test sizes $N$ per problem so that B1's
total inference FLOPs match B4's on that same problem
(Section~\ref{sec:exp-frontier}).
\textcolor{red}{\textbf{[PRE-REG TODO: confirm the single B1 FLOP-match
target is per-problem match to B4 (not the B2--B4 average); suggested:
per-problem match to B4, with B1 also reported as a full swept curve,
confirm]}}
The FLOP match is what makes the comparison honest. A method that beats
best-of-$N$ at strictly more compute has not necessarily learned
anything; a method that beats best-of-$N$ at matched compute has.

B2 isolates the contribution of the trained heads independently of
search. It evaluates fast mode (Section~\ref{sec:method-inference}):
the $Q$-head argmax is acted on directly, with no MCTS at inference.
The question B2 answers is whether the policy distilled from
MCTS-improved Bellman targets does any better than greedy decoding
from the same base model. If B2 fails to beat B0, the distillation has
not transferred the search-time improvement into the parametric heads,
and the rest of the grid loses interpretability. B3 and B4 isolate the
contribution of test-time search on top of the trained heads, at two
different simulation budgets. B3 is light search (a small number of
MCTS simulations per decision); B4 is heavy search (the largest budget
we are willing to spend per query, subject to per-problem latency
constraints). The B3-versus-B4 spread is the test-time-compute knob in
action: more search should produce monotonically better accuracy on
average, with diminishing returns as $B$ grows. Finally, B5 is the
capability ceiling. It runs a larger reasoning model (a public
distillation of an o1-style or R1-style system, for instance) under
greedy decoding, and exists to bound the achievable accuracy on the
benchmark independent of our pipeline. B5 is a reference, not a target:
beating B5 is not required for the framework's claim, but the distance
between B4 and B5 calibrates how much of the headline gap our loop has
closed.

\subsection{Compute Accounting}
\label{sec:exp-compute}

The compute axis is the single most consequential design choice in the
plan, and the one most prone to silent inflation of headline results.
We pre-commit to three metrics, reported together for every method:
\emph{(i)} total inference FLOPs per problem, summed across every LLM
forward executed at decoding time including all forwards inside any
MCTS rollout; \emph{(ii)} total LLM forward passes per problem, an
implementation-independent token count; and \emph{(iii)} wall-clock
seconds at fixed hardware (one H100), with batch size, KV-cache reuse
configuration, and decoding library version logged alongside.
Reporting all three together exposes any disagreement between them,
which would itself be evidence of a measurement bug.

We pre-commit to total inference FLOPs as the primary metric for the
headline compute-versus-accuracy frontier. The justification is
twofold. First, FLOPs are the metric used by the test-time-compute
scaling literature \citep{snell2024scaling}, so the resulting frontier
is directly comparable to prior work without unit conversion. Second,
FLOPs are the metric most resistant to post-hoc cherry-picking once
fixed: wall-clock measurements vary with batch shape and runtime
implementation in ways that invite favorable hardware choices, and raw
forward-pass counts hide the per-pass cost differences between a
$1.5$B-parameter base model and a $7$B reasoning ceiling. Pre-committing
to FLOPs and then reporting wall clock and forward-pass counts
alongside makes the primary comparison stable while preserving enough
auxiliary data to surface any measurement anomalies. Section~\ref{sec:exp-threats}
discusses compute leaks: the failure mode in which a method's
inference cost is undercounted, typically by neglecting chain-of-thought
generation tokens or the LLM forwards executed inside MCTS simulations.
Pre-registering the accounting convention is the first defense against
that failure mode.

\subsection{Headline Frontier}
\label{sec:exp-frontier}

The headline figure of the eventual empirical study is a
compute-versus-accuracy frontier; we describe it in detail here because
its axes, curves, and falsification criterion are pre-registered, but
the figure itself appears in the follow-up empirical paper, not in this
one. The $x$-axis is total inference FLOPs
per problem (from Section~\ref{sec:exp-compute}), and the $y$-axis is
accuracy as measured by the top-1 verifier pass rate. Each method in
the baseline grid contributes one curve, parameterized by the natural
compute knob for that method: best-of-$N$ varies $N$, the MCTS modes
vary $B$, and the base-greedy and ceiling baselines contribute single
points. The frontier as a whole describes how accuracy scales with
inference compute across the grid.

The pre-registered falsifiable hypothesis specifies operational
parameters in advance so that the test cannot be re-interpreted after
the empirical work is in hand. We fix the MCTS simulation budgets at
$N_{\mathrm{small}} = 50$ for the light-search baseline B3 and
$N_{\mathrm{large}} = 500$ for the heavy-search baseline B4. These
values are committed before training. They are chosen so that
$N_{\mathrm{small}}$ exceeds the verb branching factor by enough
to give every action multiple visits on average from the root, and so
that $N_{\mathrm{large}} = 10 \cdot N_{\mathrm{small}}$ tests a
one-order-of-magnitude compute scaling. Best-of-$N$ for B1 is sized
problem-by-problem so that B1's total inference FLOPs match the FLOPs
spent by B4 on that same problem; the FLOP match is the
control that makes B1 a fair cheap-deliberation baseline rather than a
straw man.

The pre-registered hypothesis is then: on GSM8K and on MATH levels $1$
through $3$, the B4 mean accuracy strictly exceeds the B1 mean accuracy
at each of $K = 5$ FLOP-matched evaluation points logarithmically
spaced across the shared FLOP range of B1 through B4, under the
statistical analysis plan of Section~\ref{sec:exp-stats}. Strict
exceedance is operationalized as a paired bootstrap $95\%$ confidence
interval on the per-problem accuracy difference excluding zero, after
Holm-Bonferroni correction across the two benchmarks and across the
$K = 5$ evaluation points (a total of $10$ tests). The hypothesis is
the strong form ``strictly above at all evaluation points'' rather
than the weak form ``above at some evaluation point'' because the weak
form invites cherry-picking of the FLOP value that reports best.
\textcolor{red}{\textbf{[PRE-REG TODO: the conjunctive ``all $10$ tests
reject'' criterion controls cherry-picking but inflates Type II risk
(diminishing returns make the top-FLOP point most likely to fail and
sink the whole claim), and no power analysis or test-set size $n$ is
stated; suggested: add a power analysis at the expected effect size, or
restrict the conjunction to the evaluation points where a gap is
expected, or pre-register a ``supported with saturation'' verdict for
failures confined to the top FLOP point, and state $n$ (the test-set
size) per benchmark, confirm]}}

If the hypothesis fails, the approximate $Q$-learning loop with
MCTS-supplied targets has not produced value beyond cheap sampling
from the same base model, and the framework's headline claim does not
hold; the appropriate response is to revise or retract the claim
rather than to add baselines or restate the result on a different
metric. The restriction to GSM8K and MATH levels $1$ through $3$
acknowledges that on AIME and MATH levels $4$ through $5$, both B1
and B4 may be near chance and the comparison is not informative; we
still report the frontier on the harder benchmarks but the
falsification criterion is attached only to the benchmarks where both
methods produce a measurable nonzero signal.

The figure also visualizes the two intermediate claims that motivate
the loop. The first is that B2 sits above B0 at the (very small)
compute of fast mode: this is the claim that distillation from
MCTS-improved targets transfers nontrivial deliberation into the
parametric $Q$-head and policy head. The second is that B3 and B4 sit
above B2 at higher compute on the same axis: this is the claim that
test-time search on top of the trained heads continues to add
accuracy. Both intermediate claims are also pre-registered as
falsifiable; failure of the first reflects a distillation failure and
failure of the second reflects either an inference-time search
implementation problem or an indication that the heads already capture
most of the local search return.

\subsection{Sanity-Check Ablations}
\label{sec:exp-ablations}

Three pre-registered ablations confirm that the trained heads, rather
than the bare MCTS scaffold, are doing the work. Each ablation
isolates one component of the trained system by replacing it with a
neutral default and re-evaluating on GSM8K and MATH levels $1$ through
$3$. Like the baseline grid, the ablation set is fixed before training
begins. The pass/fail verdicts below are stated qualitatively
(``visibly below'', ``indistinguishable from B4'', ``zero or
negative''); to remove the post-hoc degree of freedom that a qualitative
threshold introduces, each should be operationalized against the same
bootstrap machinery the headline test uses.
\textcolor{red}{\textbf{[PRE-REG TODO: commit numerical pass/fail
thresholds for A1/A2/A3 before training; suggested: reuse the paired
bootstrap CI-excludes-zero rule and the $|\bar{\Delta}| = 0.02$
effect-size band already defined in Section~\ref{sec:exp-stats}, applied
to the relevant A-vs-B4 or A-vs-B0 accuracy difference, confirm]}}

\textit{A1: untrained heads with heavy MCTS}. We replace $\Qpsi$ with a
random initialization and $\ptheta$ with the uniform distribution over
available actions, then run MCTS at $B = N_{\mathrm{large}}$ at
inference. The expectation is that A1 lies visibly below B4, and indeed
not much above B0, because PUCT with a random value head and a uniform
prior is doing nothing more than uniform exploration biased by an
uninformative leaf evaluator. If A1 matches or beats B4, then the bare
MCTS-plus-base-LLM scaffold is the source of the improvement and the
trained heads are not contributing; the framework's claim about
approximate $Q$-learning over the context MDP would not hold. A2 and
A3 are then redundant.

\textit{A2: trained heads with uniform PUCT prior}. We keep the trained
$\Qpsi$ as the leaf evaluator and the value term in PUCT, but replace
$\ptheta$ with the uniform distribution over available actions in the
PUCT exploration term. The MCTS simulation budget is matched to B4.
This ablation tests whether the distilled policy head matters beyond
what fast-mode $Q$-argmax already extracts: a policy head that
contributes nothing further than its value-head counterpart would leave
A2 indistinguishable from B4. A2 below B4 means the policy distillation
of Section~\ref{sec:method-training} is doing real work in biasing
search; A2 equal to B4 means the policy head is decorative and the
$Q$-head alone carries the contribution.

\textit{A3: trained heads, no search}. This is exactly the fast-mode
baseline B2 in the grid, but reported again here as an ablation against
B4 to make the test-time-compute claim explicit. The B4-minus-A3 gap is
the marginal value of test-time MCTS over the distilled policy alone.
If the gap is zero or negative, then test-time search adds nothing on
top of fast mode, and the inference-mode story of
Section~\ref{sec:method-inference} reduces to a single mode rather than
a knob. We pre-register A3 separately from B2 because the comparison
framing is different: B2 is a baseline in the headline frontier; A3 is
specifically the ``MCTS contributes nothing'' null result against which
the test-time-compute knob is tested.

\subsection{Statistical Analysis Plan}
\label{sec:exp-stats}

The headline hypothesis of Section~\ref{sec:exp-frontier} is tested
under a fixed statistical procedure that is committed before any
training begins. The pre-registration of the procedure is as
load-bearing as the pre-registration of the hypothesis itself: an
unspecified procedure invites the same post-hoc tuning that the
baseline grid is designed to prevent.

\paragraph{Test statistic.}
For each FLOP-matched evaluation point and each benchmark, we form the
per-problem accuracy difference $\Delta_i = \mathrm{acc}_{B4, i} -
\mathrm{acc}_{B1, i}$ across the $n$ problems in the benchmark.
Accuracy on a single problem is binary (verifier passes or fails). The
test statistic is the sample mean $\bar{\Delta} = \frac{1}{n}
\sum_{i=1}^{n} \Delta_i$.

\paragraph{Inference procedure.}
We use a paired bootstrap on the per-problem $\Delta_i$ with $R =
10{,}000$ resamples. The $95\%$ confidence interval is the
$2.5$th and $97.5$th percentile of the resampled $\bar{\Delta}$
distribution. The null hypothesis at each FLOP-matched evaluation
point is $\bar{\Delta} = 0$ at that point; rejection means the CI
excludes zero.

\paragraph{Multiple comparisons.}
The hypothesis is tested across two benchmarks (GSM8K, MATH levels
$1$ through $3$) and $K = 5$ FLOP-matched evaluation points per
benchmark, for a total of $10$ tests. We apply Holm-Bonferroni
correction at family-wise error rate $\alpha = 0.05$. The hypothesis
is supported only if all $10$ tests reject after correction.
\textcolor{red}{\textbf{[PRE-REG TODO: state the per-benchmark test-set
size $n$ used in the bootstrap, and accompany the all-$10$-reject rule
with a stated power target (or the ``supported with saturation''
fallback noted in Section~\ref{sec:exp-frontier}); confirm]}}

\paragraph{Seed variance.}
All trained methods (B2, B3, B4 and the ablations) are trained from
at least three independent random seeds. We report mean and standard
deviation across seeds for every accuracy point. The per-problem
accuracy used in the bootstrap is the mean across seeds; the
confidence interval therefore reflects problem-level variance rather
than seed variance, and seed variance is reported alongside as a
secondary diagnostic.

\paragraph{Effect size.}
Beyond the binary reject/no-reject outcome of the bootstrap, we report
the raw $\bar{\Delta}$ and its bootstrap CI at every FLOP-matched
point. A null result (failure to reject) with $|\bar{\Delta}| < 0.02$
across all FLOP points is considered evidence against the hypothesis:
the framework has neither outperformed the cheap-deliberation baseline
nor produced enough effect to suggest a power problem. A null result
with $|\bar{\Delta}| > 0.02$ but a CI that includes zero is reported
as inconclusive and prompts a power-analysis follow-up rather than a
clean accept or reject.

\paragraph{Reporting requirements.}
We commit to releasing, alongside any results paper, the raw
per-problem accuracy logs for every method and every seed, the FLOP
accounting per problem and method, the random seed values, and the
exact bootstrap implementation used. Reviewers will be able to
reproduce the bootstrap from the released artifacts; the published
test outcome will be verifiable rather than asserted.

\subsection{Threats to Validity}
\label{sec:exp-threats}

The experimental design is sensitive to several measurement and
attribution failures that have inflated headline results in prior
test-time-compute work, and we list them here as part of the
pre-registered plan rather than as discussion after the fact. Each
threat is paired with the mitigation built into the plan.

\textit{Compute leaks}. The most common failure in test-time-compute
scaling claims is undercounting inference cost. Two specific leak
modes apply to this work. The first is chain-of-thought leakage: a
method that uses token-heavy generation per decision but is scored
only on the small model's FLOPs across decoder calls undercounts the
true cost by an order of magnitude. The compute accounting of
Section~\ref{sec:exp-compute} addresses this by requiring total
inference FLOPs summed over every LLM forward, including
verb-instantiator continuations that produce action content. The
second leak mode is MCTS-rollout leakage: a method that reports only
root-level forwards but spends many forwards inside the search
procedure misrepresents the per-problem cost. The same accounting
convention catches this if applied consistently. We also report the
three compute metrics together so that any pairwise inconsistency
between FLOPs, forward counts, and wall clock surfaces as evidence of
a measurement bug rather than as a silent inflation.

\textit{Verifier coverage}. A buggy or lenient verifier inflates
accuracy across all methods, but not uniformly: methods that produce
outputs in formats the verifier handles well are advantaged over
methods whose outputs require more aggressive canonicalization. On
GSM8K and MATH this risk is concentrated in the SymPy normalization
step: a malformed regex extraction can let a numerically wrong answer
parse as correct, and a too-aggressive canonicalization can collapse
distinct symbolic answers into a single equivalence class. We will
audit a stratified sample of $200$ verifier verdicts per benchmark
before the headline run, with the audit prompts and results released
alongside the experimental code.
\textcolor{red}{\textbf{[PRE-REG TODO: commit a verifier-audit
acceptance threshold (the audited verdict error rate above which the
benchmark is flagged); suggested: mirror the $5\%$ contamination
convention used below, i.e. flag any benchmark whose audited
verifier-error rate exceeds $5\%$, confirm]}}
The verifier interface is held
constant across the baseline grid, so this risk does not preferentially
inflate the proposed method, but a buggy verifier still pollutes the
absolute numbers.

\textit{Benchmark contamination}. Pretraining corpora for the
base-model substrates may include test problems from GSM8K, MATH, and
AIME, in which case the apparent accuracy of B0 partially reflects
memorization rather than reasoning. We report a per-benchmark
contamination check (substring matching the canonicalized problem
statement against the base model's tokenizer and against publicly
documented training sets) and flag any benchmark for which the
contamination rate exceeds a pre-registered threshold of $5\%$. This
does not invalidate the framework's contribution (the trained heads
still have to learn to exploit memorized problems, and the
contamination affects all baselines symmetrically), but it does affect
the interpretation of absolute scores against frontier reasoning models.

\textit{Best-of-$N$ selector choice}. B1 requires a selector to pick a
single output from $N$ samples. The selector affects the baseline's
strength: an oracle verifier as selector trivializes B1 by reducing
best-of-$N$ to a coverage check, while a weak heuristic selector
understates B1's true accuracy. We pre-register the selector as the
same terminal verifier used to score all methods in the grid: if the
verifier accepts at least one of the $N$ samples, B1 scores it as
correct; otherwise, B1 scores it as incorrect. This is the strongest
honest baseline for B1, and it is the standard against which the B4
curve has to win.

\textit{Seed variance}. Approximate $Q$-learning with MCTS-supplied
targets is sensitive to training-seed variance, both because MCTS
trajectory collection is stochastic and because the head initialization
interacts with the early Bellman targets in ways that amplify across
the iterations. We pre-register at least three training seeds per
configuration and report mean accuracy with standard deviation across
seeds on every figure and table. Single-seed numbers are not
publishable as headline results under this plan. Inference-time
randomness (MCTS rollout sampling, temperature on the verb-instantiator)
is averaged within seed: for each test problem we run at least eight
inference replicates per method, with seed-specific aggregation
described in the accompanying release notes.

\textit{Selection effects in the pre-registered hypothesis}. The
falsification criterion in Section~\ref{sec:exp-frontier} restricts
attention to GSM8K and MATH levels $1$ through $3$ because both B1 and
B4 produce measurable signal there. This restriction is itself a
pre-registered methodological choice rather than a degree of freedom
discovered after seeing results: the restriction is fixed before any
empirical work, and the harder-benchmark frontiers are reported
without a falsification claim attached. The published study will state
which benchmarks contribute to the falsification criterion and which
contribute only descriptively, in advance of any number-reporting.
